\section{Pereshi — ``The Roof of the Way''}
\label{chap:pereshi}

\subsection*{Quick Reference}
\textbf{Tagline:} \textit{The fire remembers. The archive endures. The Fhara fear us; the Tulkani remember us; the lowlands do not see us. That is how we like it.}

\begin{quote}
  \textbf{Flame-Keeper (Elite):}  
  ``The fire does not forget. It has witnessed every oath sworn in this valley, every treaty sealed with a dropped coal, every lie told in its light. I do not tend the flame. I listen to what it remembers. And it remembers everything – including the Prince who thought he could burn us. We took his name. We cut his shadow. But shadows grow back, and the ninth coal is warm.''
\end{quote}

\begin{quote}
  \textbf{Caravan Scribe (Commoner):}  
  ``My grandfather copied a tax ledger that had not been opened in three centuries. He found a mistake – a single miscounted bale of silk. That mistake started a war. We do not write for the living. We write for the dead who cannot correct their errors. And we write to remind ourselves: the empire of ash will not rise again. But the Prince's shadow is long, and the ninth coal has started to hum. The Fhara send envoys who cannot meet our eyes. The Tulkani leave knots at our thresholds. We do not explain. We burn.''
\end{quote}

\textbf{Genre / Mood:} High alpine bureaucracy, firelit memory, postcolonial witness, cursed return.

\textbf{Starting Location:} The scriptorium of a mountain caravanserai, where lamps burn at noon and scribes copy tax ledgers by candlelight. A Flame-Keeper looks up from a brazier: ``You have come for a tariff receipt, or for a verse to cross the pass, or to consult a book that should have been burned. First, tell me a truth – something the Ashaani Prince tried to erase. Then I will decide which fire shows you the way. And be quick: the ninth coal is warmer than it should be. It has been humming in my sleep.''

\vspace{2mm}
\begin{tcolorbox}[colback=black!3,colframe=blue!60!black,title={Theme \& Atmosphere}]
The Pereshi highlands are a vast, windswept plateau, cut by deep river valleys and crowned with snow-capped peaks. Ancient fire temples stand beside caravanserais; terraced vineyards climb the hillsides. The people are scholars, shepherds, and traders – proud, suspicious, and bound by a code of hospitality that demands three days of feasting for any guest, and three generations of vigilance for any insult.

But the Pereshi are not a nation. They are a memory of one. Long ago, a highland empire stretched from the western passes to the eastern salt flats – a confederation of fire temples, fortress monasteries, and clan valleys that answered to a council of Flame‑Keepers. That empire fell when the Ashaani Prince marched east, burning libraries and breaking the hearths. What remains is a scatter of independent holdings – cities like Mount Khvarena, temple towns like Adur‑Gah, and countless villages that recognise no king, only the fire and the ninth coal in every hearth.

The Pereshi share a culture: the same fire rituals, the same oral epic (the \textit{Book of Ash}), the same fear of naming the dead. But they do not unite. They have no crown, no single council. The Flame‑Keepers meet once every nine years to renew the seal on the Prince’s heart, but they do not govern. They witness. The rest of the time, each valley minds its own thresholds. This fragmentation is a shield – it has kept conquerors from finding a single throat to cut. It has also kept the Pereshi from reclaiming what they lost.

\vspace{2mm}
\textbf{What you notice first:} The smell of woodsmoke and old parchment. The way every valley has a different accent. The constant call to prayer from fire‑temple towers – but the prayers are in a language no one speaks outside the highlands, a language that predates the Prince's rule. And the way every Flame‑Keeper touches the ninth coal before sleep, checking if it has grown hotter. Lately, they check twice.

\textbf{The one rule that kills newcomers:} Never step on a fallen book. Written words carry the weight of law – and the weight of witness. To step on a book is to trample a contract, a genealogy, a verdict, a memory of the Prince's cruelty. The spirits of the words will haunt you. Also, never speak the Prince's true name aloud. The scroll remembers. The shadow will hear. And the ninth coal will pulse.
\end{tcolorbox}

\subsection*{Regional Mood: The Weight of the Archive and the Warmth of the Coal}

Power here is measured in fire‑temple charters, caravan tariffs, and the length of a family's library. The petty lords of the highlands compete not through armies but through scribes who can produce the oldest charter, the most complete genealogy, the most beautifully illuminated tax receipt. The Covenant has a strong presence in the sacred cities, but the Pereshi follow their own law: the word of a Flame‑Keeper outweighs a judge's gavel. Travelers are expected to recite a verse when crossing a bridge – the spirits of the water demand it – and to present a red rose before petitioning a lord. The high passes are measured in breath‑counts, not miles, and the air is thin enough to make a foreigner's head spin. But the thinnest air of all is in the sealed archive where the Prince's name is buried. No one has opened it in nine generations. The lock has no key. The key is a memory no one wants to have – and lately, the lock has been weeping rust. The rust tastes of iron and old smoke.

\subsection*{The Unfired Crown: How the Pereshi Rule Themselves}
\label{subsec:pereshi-unfired-crown}

The Pereshi do not kneel to a king. They kneel to fire, to ash, and to a name they erased. Each valley or fire‑temple city governs itself through a council of elders, a Flame‑Keeper, and a war captain (if needed). There is no central treasury, no standing army, no code of laws beyond custom and the witness of the ninth coal. What unites them is not a crown but a memory of conquest: every Pereshi child learns the story of the Ashaani Prince, the burning of the great library at Khvarena, and the year the fires went out. They also learn the year the fires were relit – and the ritual that severed the Prince's name from the world.

The Flame‑Keepers of the nine largest temples meet once every nine years at Mount Khvarena, the holiest peak. They do not elect a leader. They sit in a circle, each holding a coal from their own hearth. They speak of weather, of harvests, of bandit raids. They do not speak of the Prince – but each of them, before sleep, places a hand on the ninth coal that sits in a brass brazier at the centre of the circle. If it is warm, they know the Shadow is stretching. If it is hot, they light a signal fire and send riders to every valley: \textit{The name stirs. Renew the seals.}

No outsider has ever attended a meeting of the nine. The Fhara have tried to send spies; the spies never returned, and their names were erased from every ledger. The Tulkani do not try; they leave a knotted cord at the edge of the temple and wait for a reply. The reply is always a single burnt knot.

\subsection*{The Fear of the Fhara and the Tulkani Unease}
\label{subsec:pereshi-fear-and-unease}

\textbf{The Fhara fear the Pereshi.} In the desert oases, water‑clerks whisper that every Pereshi carries a ``name‑knife'' – a coal that can erase a debtor from memory. A merchant whose name is struck from the ledgers is not bankrupt; he is \textit{forgotten}. His wells dry up. His caravans lose their way. His children cannot remember his face. The Fhara pay their debts to the Pereshi with unusual haste, and they never, ever allow a Pereshi fire‑keeper to light a lamp in their counting houses. Some say that the ninth well of Qalat al‑Madina was cursed not by the Kuvani but by a Pereshi exile who wanted to make an oasis vanish.

\textbf{The Tulkani have an uneasy relationship with the Pereshi.} They share a border of silence. Both respect the ninth night, but for different reasons. The Tulkani say the Pereshi stole fire from the Dreamer (Ikasha). The Pereshi say the Tulkani forgot how to build a hearth. They trade stories at the edge of the high passes, but never stay the night. A Tulkani who enters a Pereshi fire temple must leave a knot behind – and the Pereshi burn it. The smoke carries the dream to the sky, but the knot is gone, and the Tulkani cannot retie it. Some Tulkani elders whisper that the Pereshi know how to sever a name because they once tried to sever the Dreamer's name – and failed. They do not say this aloud. The Hollow might be listening.

\textbf{Everyone else underestimates the Pereshi.} Ecktorian clerks call them ``backward highlanders with no coin.'' Vhasian knights sneer at their unpaved roads and unarmoured sentinels. Oshiiran grain factors pity them. The Pereshi prefer it this way. Let the lowlands send their tax collectors; the passes are high, the cold is deep, and every tariff receipt requires a verse that no foreigner can recite correctly. Those who try are turned back at the bridge – or find that their names have been quietly written in ash and left to the wind.

\subsection*{The Forgotten Prince – Heart as Coal, Name as Ash}
\label{subsec:pereshi-forgotten-prince}

Once, the Ashaani Prince – a conqueror from the Veiled Throne, a man whose armies swept from the Red Desert to the eastern hills – turned his gaze to the Pereshi highlands. He burned the great library at Khvarena. He melted the temple bells into arrowheads. He tried to erase the fire‑faith, to replace it with the silence of the Veiled Sisters. For nine years he ruled from a garrison in the lowlands, sending patrols to cut down any who still lit a lamp before dawn.

Then his empire cracked. The Kalrantha fragment that had given him power over fate began to devour him from within. His generals turned on each other. His shadow began to walk independently. He retreated to a fortress in the eastern passes, and the Pereshi followed – not with armies, but with priests.

The ritual took nine nights. They did not kill him. They \textit{unmade} him. They spoke his true name into a lead scroll, sealed it with nine coals, and buried it beneath the largest fire temple. They cut his shadow from his body and pinned it to the Sun King's Throne (now lost in the Red Desert). And they took his heart – still beating, still warm – and placed it in a brass brazier. The ninth coal is not a relic. It is his petrified heart. Once a year, on the anniversary of his fall, it pulses. The Flame‑Keepers feel it. They do not speak of it.

The Prince's shadow still walks because the heart still dreams of conquest. The Pereshi do not fear him – they \textit{remember} him. He is the Debt That Was Paid. But the debt compounds. Every broken name, every forgotten oath, every coal that cracks adds interest. If the ninth coal ever goes cold, the Prince will not return – he will never have been gone.

\subsection*{The Prince's Shadow Clock [8]}
Tracks the Prince's gradual return. Advance the clock when:
\begin{itemize}
  \item The party breaks a Pereshi oath or defiles a fire temple (+1)
  \item A Flame‑Keeper dies without passing on the secret of the ninth coal (+1)
  * The party speaks the Prince's true name aloud (from the First Emperor's scroll) (+2)
  \item The GM spends 2 SB to represent a crack in the archive's seal (+1)
\end{itemize}
When the clock reaches 4: The Prince's shadow begins to appear in mirrors and still water. It asks questions. Answering gives it a thread to pull. When the clock reaches 8: The Prince manifests as a ghostly figure – not a combat boss, but a presence that can possess a willing (or desperate) host. The party must either reseal him (by finding a new name, a new shadow, or a new coal) or bargain with him (he offers great power at the cost of a Pereshi valley). The valleys are not his to give. He offers them anyway. The ninth coal will be warm that night.

\subsection*{Prominent NPCs of the Pereshi Highlands}
\label{sec:pereshi-npcs}

\begin{longtable}{|p{0.2\textwidth}|p{0.25\textwidth}|p{0.3\textwidth}|p{0.15\textwidth}|}
\hline
\textbf{Name} & \textbf{Role/Title} & \textbf{Motivation} & \textbf{Rumoured Quirk} \\
\hline
Flame‑Keeper Zarifa & Priestess of Mount Khvarena & Preserve the old rites and keep the ninth coal from lighting & She has never been seen without a rose – fresh or dried. The rose is from the Prince's own garden, taken the day he fell. She sleeps with one hand on the coal. The coal has begun to pulse in time with her heart. \\
\hline
Lord Rostam & Highland petty king (Adur‑Gah valley) & Unite the valleys against lowland tax collectors – and against the Prince's return & He keeps a cage of doves; each dove carries a different treaty. The ninth dove carries a blank page for the Prince's name. The dove has started to speak. It speaks in the Prince's voice. \\
\hline
Archivist Kaelen & Keeper of the Unfinished Books (Khvarena) & Protect the scrolls that were never completed – especially the Prince's name scroll & He has not left the archive in thirty years. His shadow is missing. The Prince's shadow took it. He does not sleep. He no longer needs to. \\
\hline
Tarkan & Caravan master of the high passes & Secure the passes before the snows close them – and before the Prince's ghost walks & He has three camels, each named after a dead wife. The camels were bred from stock the Prince's army abandoned. They refuse to cross the Pass of the Unforgiven. They kneel when the ninth coal pulses. \\
\hline
The Poet of the Ninth Verse & An old man who stole a missing line of an ancient epic & Recite the verse for a cup of tea made with snowmelt – and keep the Prince from hearing it & His left eye is a piece of amber with a fly trapped inside. The fly is from the Prince's own stable. When the Prince's shadow is near, the fly buzzes. When the coal pulses, the fly is silent. \\
\hline
Fire‑Keeper Zahra (new) & Young priestess of the eastern temple (Adur‑Gah) & Secretly correspond with a Fhara water‑clerk – she is torn between duty and mercy & She writes letters on bark. The bark is from a grove that the Prince burned. She signs with a drop of her own blood. The Fhara clerk has never seen her face. \\
\hline
The Knot‑less (Tulkani exile, new) & A Tulkani who cut her own cord, now living among the Pereshi & Learn the name‑severing ritual – to hide from the Hollow & She carries a single knot in her mouth. She has never swallowed it. She speaks by moving her lips around the knot. The Pereshi have not yet burned her offerings. \\
\hline
The Nameless (Fhara merchant, new) & A spice trader whose name was erased by a Pereshi fire‑keeper & Reclaim his lineage – even if it costs a shard of the Kalrantha amulet & He wears a mask of woven silver. Underneath, his face is smooth – no mouth, no nose, only two eyeholes. He drinks through a straw. He whispers his true name to the ninth coal every night, hoping it will remember. \\
\hline
\end{longtable}

\subsection*{Names of the Pereshi Highlands}
\label{sec:pereshi-names}

Inspired by highland and steppe traditions: names with soft endings, honorific prefixes (Mir, Khan, Beg), and surnames denoting a clan or a valley. Many families carry a second, secret name – the name they used during the Prince's occupation, preserved as a mark of survival. To speak a Pereshi’s secret name is to bind them to a promise; to learn it is to hold a thread of their soul.

\begin{longtable}{|p{0.25\textwidth}|p{0.25\textwidth}|p{0.2\textwidth}|p{0.2\textwidth}|}
\hline
\textbf{First Names (Male)} & \textbf{First Names (Female)} & \textbf{Surnames} & \textbf{Epithets} \\
\hline
Rostam & Zarifa & Kuhistani & \textit{the Fire‑Touched} \\
\hline
Kaelen & Laleh & Qasemi & \textit{the Unfinished} \\
\hline
Tarkhan & Shirin & Salari & \textit{Ninth Verse} \\
\hline
Dariush & Mahsa & Bahrami & \textit{the Bridge‑Mouthed} \\
\hline
Behrouz & Golnar & Khorasani & \textit{the Red Rose} \\
\hline
\end{longtable}

\paragraph*{(Temple/Library/Bridge)} Fire temple with a copper dome and a flame that never goes out; library of unfinished books where scribes copy texts that were burned by the Prince's order; stone bridge over a chasm, each arch carved with a line of poetry – except the ninth, which is deliberately blank. The blank arch hums at midnight.

\section*{Spades — Places}
\label{sec:pereshi-places}
\begin{enumerate}
\setcounter{enumi}{1}
\item \textbf{Fire Temple of the Unquenched} – A stone tower with a copper dome. The fire inside is said to be the first flame. It has never gone out – but it flickers when a king dies.
  \hfill \textit{The flame's colour changes with the season. The ninth coal sits in a brass brazier, untouched since the Prince's fall. Recently, the coal has started to glow at midnight. The glow casts no shadow.} `[RITUAL]` `[MEMORY]`
\item \textbf{Mount Khvarena (The Nine Fires)} – The holiest peak, where the Flame‑Keepers meet once every nine years.
  \hfill \textit{At the summit, a circle of nine braziers. The ninth brazier holds the Prince's heart. The coals never cool. Pilgrims climb the mountain barefoot, carrying a single rose. The rose is burned. The ash is the offering.} `[LEGEND]`
\item \textbf{The Archive of Severed Names} – A cave behind a waterfall, where clay tablets record every name ever erased.
  \hfill \textit{The tablets are stacked to the ceiling. Some are warm. Those names are still trying to return. The archivist has never opened the ninth row.} `[DEATH]` `[SECRET]`
\item \textbf{Adur‑Gah (Fire Temple City)} – A walled town built around a natural gas vent that burns day and night.
  \hfill \textit{The vent is called the Prince's Breath. The locals say it is his sigh, trapped underground. When the wind shifts, the flame roars – and the ninth coal pulses.} `[COMMUNITY]`
\item \textbf{The Pass of the Unforgiven} – A high defile where the Prince's army froze to the rock.
  \hfill \textit{On cold nights, you can still see their shapes in the ice. They are still standing. Their lips move. They are still counting. The Pereshi do not use this pass. They say the dead do not forgive.} `[DEATH]` `[FEAR]`
\item \textbf{The Library of Unfinished Books} – A scriptorium where scribes copy texts that the Prince tried to burn.
  \hfill \textit{The books have no endings. The endings were lost. The scribes write new ones, but the old endings whisper from the margins. The ninth shelf is empty. It is reserved for the book that will contain the Prince's true name.} `[KNOWLEDGE]`
\item \textbf{The Bridge of Nine Echoes} – A rope bridge over a frozen chasm. Each step echoes nine times.
  \hfill \textit{The ninth echo is always a whisper – the name of the next person to fall. Do not answer. The whisper is not for you.} `[DEATH]` `[OMEN]`
\item \textbf{The Rose Garden of Khvarena (Ruined)} – Where the Prince cut every rose bush to build a pyre.
  \hfill \textit{One bush survived. It grows from a crack in the stone. The roses are black. They smell of ash. A single petal can bind a name – or sever it.} `[MAGIC]`
\item \textbf{The Well of Unspoken Names} – A dry well where the Pereshi throw coals that have gone cold.
  \hfill \textit{The coals burn again at the bottom. No one knows how. Some say the Prince's heart keeps them warm.} `[MYSTERY]`
\item[J] \textbf{The Sunken Archive (Fhara border)} – A flooded chamber where water‑tablets and fire‑tablets meet.
  \hfill \textit{The water is warm. The tablets float. The Fhara water‑clerks who come here never leave. The Pereshi do not explain why.} `[THRESHOLD]`
\item[Q] \textbf{The Silent Pass} – A high mountain pass where no bird flies.
  \hfill \textit{The air is too thin for sound. Arguments must be written in the snow. The snow does not lie. The Pereshi leave offerings of silence here.} `[WEATHER]`
\item[K] \textbf{The Ninth Coal’s Vault} – A lead‑lined chamber beneath Mount Khvarena.
  \hfill \textit{The door has nine locks. The keys are held by nine Flame‑Keepers. The ninth key is a memory that no one can recall. The door sweats.} `[SEAL]`
\item[A] \textbf{The First Emperor’s Scroll (Myth)} – A scroll that contains the Prince's true name, written in ash on lead.
  \hfill \textit{It is kept in a box that has no key. The box is opened only when the ninth coal cracks. The last time it cracked, the Prince's shadow walked for nine days.} `[LEGEND]`
\end{enumerate}

\paragraph*{(Priest/Scholar/Envoy)} Flame‑Keeper with a rose and a coal; archivist with ink‑stained fingers; Fhara merchant with a woven silver mask.

\section*{Hearts — People \& Factions}
\label{sec:pereshi-people}
\begin{enumerate}
\setcounter{enumi}{1}
\item \textbf{Flame‑Keeper Zarifa} – Priestess of Mount Khvarena, keeper of the ninth coal.
  \hfill \textit{She has never been seen without a rose – fresh or dried. The rose is from the Prince's own garden. She sleeps with one hand on the coal. The coal has begun to pulse in time with her heart.} `[RITUAL]`
\item \textbf{Archivist Kaelen} – Keeper of the Unfinished Books, guardian of the Prince's name scroll.
  \hfill \textit{He has not left the archive in thirty years. His shadow is missing. The Prince's shadow took it. He does not sleep. He no longer needs to.} `[SECRET]`
\item \textbf{Caravan Master Tarkan} – A highland trader who knows every pass and every ghost.
  \hfill \textit{He has three camels, each named after a dead wife. The camels were bred from stock the Prince's army abandoned. They refuse to cross the Pass of the Unforgiven. They kneel when the ninth coal pulses.} `[TRAVEL]`
\item \textbf{Lord Rostam} – A petty king of Adur‑Gah, trying to unite the valleys.
  \hfill \textit{He keeps a cage of doves; each dove carries a different treaty. The ninth dove carries a blank page for the Prince's name. The dove has started to speak. It speaks in the Prince's voice.} `[POLITICS]`
\item \textbf{The Poet of the Ninth Verse} – An old man who stole a missing line of the epic.
  \hfill \textit{His left eye is a piece of amber with a fly trapped inside. The fly is from the Prince's own stable. When the Prince's shadow is near, the fly buzzes. When the coal pulses, the fly is silent.} `[KNOWLEDGE]`
\item \textbf{Fire‑Keeper Zahra (new)} – A young priestess of the eastern temple, secretly writing to a Fhara water‑clerk.
  \hfill \textit{She writes letters on bark. The bark is from a grove that the Prince burned. She signs with a drop of her own blood. The Fhara clerk has never seen her face.} `[CONFLICT]`
\item \textbf{The Knot‑less (Tulkani exile, new)} – A Tulkani who cut her own cord, now living among the Pereshi.
  \hfill \textit{She carries a single knot in her mouth. She has never swallowed it. She speaks by moving her lips around the knot. The Pereshi have not yet burned her offerings.} `[MYSTERY]`
\item \textbf{The Nameless (Fhara merchant, new)} – A spice trader whose name was erased by a Pereshi fire‑keeper.
  \hfill \textit{He wears a mask of woven silver. Underneath, his face is smooth – no mouth, no nose, only two eyeholes. He drinks through a straw. He whispers his true name to the ninth coal every night, hoping it will remember.} `[TRAGEDY]`
\item \textbf{Covenant Judge} – A grey‑robed notary stationed at the bridge between Pereshi and Vhasian lands.
  \hfill \textit{She carries a water‑clock that runs on melted snow. It never freezes. She has never ruled against a Pereshi – not because she fears them, but because their contracts are written in ash, and ash has no appeal.} `[JUSTICE]`
\item[J] \textbf{The First Censor (Ghost of Khvarena)} – A librarian who burned herself with the scrolls to keep the Prince from reading them.
  \hfill \textit{She appears as a pillar of smoke. She asks: ``Is the name still sealed?'' If you say yes, she weeps ash. If you say no, she screams and the ninth coal cracks.} `[DEATH]`
\item[Q] \textbf{The Prince's Shadow (Manifestation)} – A silhouette without a face, seen in mirrors and still water.
  \hfill \textit{It asks for a name. Any name. Give it a living one, and that person will dream of the Prince every night. Give it a dead one, and the dead will walk for one night. Give it your own, and you lose your shadow – and your name becomes ash.} `[CURSE]`
\item[K] \textbf{The Ninth Flame‑Keeper (Legend)} – The one who sealed the Prince's heart. She never died.
  \hfill \textit{She sits at the bottom of the Well of Unspoken Names, warming the coals with her breath. She has not breathed in nine centuries. The coals breathe for her.} `[MYTH]`
\item[A] \textbf{The Prince's Echo (Forgotten Conqueror)} – A voice from the ninth coal, heard when it pulses.
  \hfill \textit{It speaks in a language no one remembers. It says only one word. The word is a name – yours, or the Prince's. It changes each time.} `[FEAR]`
\end{enumerate}

\paragraph*{(Ash/Name/Coal)} A tablet from the Archive of Severed Names (contains an erased name); a rose from the Ruined Garden (can be burned to ask a yes/no question of the fire); a fragment of the ninth coal (once, reroll a failed negotiation with a Pereshi).

\section*{Clubs — Complications/Threats}
\label{sec:pereshi-complications}
\begin{enumerate}
\setcounter{enumi}{1}
\item \textbf{Name‑Debt Called} – A Pereshi fire‑keeper demands a name you do not wish to give.
  \hfill \textit{The name must be true. It must be someone you know. The fire‑keeper will write it on a clay tablet and place it in the Archive of Severed Names. That person will begin to be forgotten – first by strangers, then by family, then by you.} `[DEBT]` `[SOCIAL]`
\item \textbf{The Ninth Coal Cracks} – The Prince's heart splits. The Shadow stretches.
  \hfill \textit{Advance the Prince's Shadow clock by 2. All mirrors in the scene show a silhouette standing behind each character. The silhouette asks: ``Why do you keep my name?'' Refuse to answer, and it will take a memory.} `[DEATH]` `[FEAR]`
\item \textbf{Fhara Envoy Refused} – A water‑clerk from the oases arrives with gifts. The Pereshi burn them.
  \hfill \textit{The envoy is the Nameless – a merchant whose name was erased. He offers the party a shard of the Kalrantha amulet in exchange for a promise to intercede. The Pereshi will not negotiate. The party must choose between the merchant's gratitude and the fire‑keeper's ire.} `[POLITICS]`
\item \textbf{Tulkani Knot Left at the Threshold} – A Tulkani Caravan Judge has left a knotted cord outside a fire temple.
  \hfill \textit{The knot is a request for parley. The Pereshi will not touch it. They will burn it at dusk. If the Tulkani returns before the fire goes out, they may speak – but only in riddles. The party is asked to act as interpreters. A wrong answer may start a feud.} `[SOCIAL]`
\item \textbf{Ash from the Prince's Pyre} – A storm carries ash from the Red Desert. The ash settles on Pereshi fields.
  \hfill \textit{The ash is warm. It smells of old smoke. The Pereshi say it is the Prince's breath, scattered. Any promise written in it will be forgotten by dawn. The party finds a message written in ash on their tent. It says: ``You carry my name.''} `[OMEN]`
\item \textbf{The First Emperor's Scroll Found} – A copy of the scroll containing the Prince's true name surfaces in a lowland market.
  \hfill \textit{The Pereshi will pay anything – a name, a coal, a memory – to recover it. The party is hired to retrieve it. But other factions want it too: the Veiled Sisters (to free the Prince), the Breath‑less (to destroy it), and a Fhara merchant (to sell it to the highest bidder).} `[QUEST]`
\item \textbf{The Ninth Verse Spoken Aloud} – Someone recites the missing line of the epic.
  \hfill \textit{The line is: ``And the heart that burned, the name that ash – the ninth is warm, the shadow waits.'' The moment it is spoken, the ninth coal pulses. A Oath‑Keeper appears. It asks: ``Who taught you that verse?'' The speaker must answer truthfully or lose their voice.} `[CURSE]`
\item \textbf{Fire Temple Silence} – All the fires in a temple go out at once. The ninth coal is cold.
  \hfill \textit{The Flame‑Keeper is found dead, a rose in her mouth. The ninth coal is not cold – it is gone. Stolen. The party must track the thief before the Prince's shadow claims the coal. The thief is a Tulkani exile who thought the coal would let her cut her last knot.} `[CRIME]`
\item \textbf{The Archive Weeps} – Water seeps from the walls of the Archive of Severed Names.
  \hfill \textit{The water tastes of salt and old ink. It is the tears of the erased names. They are trying to return. A character who touches the water hears a single name whispered – a name they have never heard. That name now belongs to them. The Hollow will expect it.} `[DEBT]`
\item[J] \textbf{The Prince's Shadow Points} – The shadow appears in a mirror and points at a party member.
  \hfill \textit{That character now owes a name to the Prince. They have nine days to find someone else's name to offer, or the Prince will take theirs. The shadow will appear in every reflection until the debt is paid.} `[CURSE]` `[DEBT]`
\item[Q] \textbf{The Rose Garden Burns (Again)} – Someone sets fire to the ruined garden. The black roses scream.
  \hfill \textit{The smoke forms the shape of a crown – the Prince's crown. The fire was lit by a Fhara merchant who thought he could destroy the Prince's memory. Instead, he woke it. The Prince's shadow now walks toward the lowlands. The party must warn the Pereshi – or bargain with the Prince.} `[DISASTER]`
\item[K] \textbf{The Ninth Flame‑Keeper Rises} – The legendary priestess climbs out of the Well of Unspoken Names.
  \hfill \textit{She has not aged. Her eyes are coals. She says: ``The seal is failing. I need a new name to bind.'' She looks at the party. She asks: ``Who volunteers?'' Refusal will not be accepted. The party must offer a false name, a dead name, or a memory of a name. The flame‑keeper will know if it is false.} `[LEGEND]` `[FEAR]`
\item[A] \textbf{The Prince Walks} – The ninth coal cracks completely. The Prince's shadow becomes solid.
  \hfill \textit{He does not fight. He negotiates. He offers: ``Give me a kingdom, and I will give you the world. Give me a name, and I will forget yours. Give me a coal, and I will warm your hearth forever.'' The Pereshi will not help. They are hiding. The party must decide alone.} `[APOCALYPSE]`
\end{enumerate}

\paragraph*{(Ash/Name/Coal)} Fragment of the ninth coal (once, reroll a failed negotiation with a Pereshi); a rose from the Ruined Garden (burn it to ask a yes/no question of the fire); a tablet from the Archive of Severed Names (contains an erased name – can be used to blackmail someone who has forgotten that name).

\section*{Diamonds — Rewards/Leverage}
\label{sec:pereshi-rewards}
\begin{enumerate}
\setcounter{enumi}{1}
\item \textbf{Fire‑Temple Token (Copper)} – A small disc stamped with a flame. Gain +1 Position in any Pereshi court or negotiation.
  \hfill \textit{The token is warm. It cools when you lie. The Pereshi will notice.} `[SOCIAL]`
\item \textbf{Rose from the Ruined Garden} – A black rose, dried, from the bush that survived the Prince's pyre.
  \hfill \textit{Burn it to ask a yes/no question of the fire. The fire will answer with a colour: red for yes, blue for no, white for the question is forbidden. The rose crumbles after use.} `[MAGIC]`
\item \textbf{Tablet of a Severed Name} – A clay tablet from the Archive, inscribed with a name that was erased.
  \hfill \textit{The name belongs to a person who has been forgotten – even by themselves. Speak the name aloud, and that person will remember who they are. Use it to blackmail someone who once owed that person a debt. The tablet crumbles after one use.} `[SECRET]` `[DEBT]`
\item \textbf{Fragment of the Ninth Coal} – A piece of the Prince's heart, still warm.
  \hfill \textit{Once, reroll a failed negotiation with a Pereshi. The fragment then turns to ash. But the Prince will remember your face. The Hollow will too.} `[LEGEND]`
\item \textbf{The Ninth Verse (Scroll)} – The missing line of the epic, written on a scrap of leather.
  \hfill \textit{Recite it aloud to summon a Oath‑Keeper. The Oath‑Keeper will answer one question about the Prince's shadow – but it will then ask you a question in return. You must answer truthfully or lose your voice.} `[KNOWLEDGE]` `[DEBT]`
\item \textbf{A Coal from Mount Khvarena} – A regular coal, blessed by the ninth flame.
  \hfill \textit{Carry it to be immune to cold weather for one leg of a journey. The coal glows faintly when the Prince's shadow is near.} `[TRAVEL]`
\item \textbf{Archivist's Seal} – A clay seal that opens any archive door in the highlands.
  \hfill \textit{The seal is shaped like a rose. It must be pressed to the lock. The lock will open – but the archivist will know who opened it. The seal works once.} `[ACCESS]`
\item \textbf{The Prince's True Name (in a lead scroll)} – A copy of the name, written in ash.
  \hfill \textit{Speak it aloud, and the Prince's shadow will manifest and obey one command. Then the scroll burns, and the Prince will hunt you. Use once. The GM gains 5 SB.} `[LEGEND]` `[CURSE]`
\item \textbf{Fhara Water‑Clerk's Ledger (Stolen)} – A ledger that records every name the Pereshi have erased for Fhara merchants.
  \hfill \textit{The ledger is a weapon. Present it to a Fhara well‑judge, and they will cancel a debt. Present it to a Pereshi fire‑keeper, and they will grant you a favour – but they will also ask whose name you want to erase in return.} `[TRADE]` `[DEBT]`
\item[J] \textbf{The Knot‑less's Last Knot} – The knot that the Tulkani exile carries in her mouth.
  \hfill \textit{If you untie it, the exile will remember her name – but she will also owe you a story. If you cut it, she will die, and you will gain the ability to hide from the Hollow for one night (the Hollow cannot see you). The knot is warm. It smells of sleep.} `[MYSTERY]`
\item[Q] \textbf{The Nameless's True Name (Whispered)} – The name that the Fhara merchant lost.
  \hfill \textit{Whisper it to him, and he will grant you a boon: a shard of the Kalrantha amulet, or a map to the archive of severed names. But the Pereshi will know you spoke his name. They will mark you.} `[TRADE]` `[DEBT]`
\item[K] \textbf{The Ninth Flame‑Keeper's Rose} – A rose that never wilts, from the legendary priestess.
  \hfill \textit{Once, present it to any Pereshi to demand a single act of mercy – the release of a prisoner, the return of a name, a night's shelter during a storm. The rose will then turn to ash. The priestess will ask why you used it. Answer poorly, and she will take a memory.} `[LEGEND]`
\item[A] \textbf{The Prince's Shadow (Bound)} – A piece of the Prince's shadow, captured in a lead‑lined mirror.
  \hfill \textit{Break the mirror, and the shadow will fight for you for one scene (Cap 3, Harm 3). Then it will turn on you. The mirror is heavy. It whispers. It knows your name.} `[MAGIC]` `[CURSE]`
\end{enumerate}

\section*{Quick use notes}
\label{sec:pereshi-quick-use}
\begin{itemize}
\item Draw until you have all four suits: Spade = place, Heart = actor, Club = pressure, Diamond = leverage. Highest rank sets the main Clock (2–5 → 4, 6–10 → 6, J/Q/K → 8, A → 10).
\item Diamonds are codified outcomes (tokens, scrolls, coals) that change position rather than call for a roll.
\item If any A appears, echo \textbf{fire‑omens} – flames that burn green, scrolls that ignite without heat, a bridge that recites a verse in a dead language, a coal that glows in the dark. Also tick the Prince's Shadow clock by 1 (the curse advances). The coal is not a metaphor. It is warm.
\end{itemize}

\subsection*{Additional Features (Cultural Mechanics)}
\label{sec:pereshi-mechanics}

\begin{itemize}
\item \textbf{Fire Temples and the Ninth Coal:} Every fire temple keeps a brazier with nine coals. The ninth coal is never lit – it is a promise to the Hollow and a ward against the Prince. Lighting it summons the Oath‑Keeper or the Prince's ghost. Flame‑Keepers inherit the secret of which coal is the ninth; it changes with each keeper. The coal is warm to the touch. It should not be. Recently, all the ninth coals across the highlands have grown warm at the same time. They pulse in the same rhythm. The rhythm is not a heartbeat.
\item \textbf{Verse Law:} In Pereshi courts, a well‑recited verse carries more weight than a witness. The old epics contain every legal precedent; a litigant who can quote the correct verse wins automatically. The problem is that the epics are incomplete – and the missing verses are hidden in the tombs of forgotten kings. The Prince tried to finish the epic with his own verses. They are still there, in the sealed archive, waiting to be burned – or to be sung. The singers who try to learn them lose their voices for a year.
\item \textbf{Caravan Tariffs as Currency:} A paid tariff receipt can be traded, sold, or used as proof of credit. The Pereshi have no central bank; the receipt is the coin. Forged receipts are a capital crime – the forger is forced to copy the entire epic by hand, blindfolded. The Prince's tax collectors were the first to be executed this way. Their ghosts still copy, still blindfolded. The ink they use is rust‑coloured.
\item \textbf{The Prince's Seal (Forgotten):} In the deepest vault of the archive, sealed in a lead box, is the Prince's own signet ring. It is said to grant authority over any Pereshi – because it proves the Prince's claim. No one has opened the box. The key was buried with the last emperor. The ring is still warm. Recently, the box has begun to hum. The hum is the same pitch as the ninth coal's pulse.
\item \textbf{Name‑Debt Mechanics:} Instead of tracking coin, a Pereshi may demand a “name” – the true name of an enemy, a forgotten ancestor, a secret that has never been spoken. A character who cannot pay loses one Bond or one Asset (the name is erased from records and memory). The GM may offer a choice: pay a name you know, or lose one of yours. This is not a debt. It is an erasure.
\item \textbf{Fire‑Temple Hospitality:} The three‑day feast rule is sacred. Breaking it marks the offender as \textit{anathema} – all Pereshi refuse to speak to them, and the ninth coal in every hearth spits ash in their direction. Removing the mark requires a pilgrimage to Mount Khvarena and the offering of a true name to the ninth coal. The coal may accept or refuse. It usually refuses.
\item \textbf{Local Patrons (Syncretic Names):}
  \begin{itemize}
    \item \textbf{The Undying Flame (Oath of Flame and Light):} The fire that witnessed the Prince's defeat – and still burns in witness and in warning. The warning is the coal's warmth.
    \item \textbf{The Keeper of Sealed Names (The Witness):} The archivist who records what must never be spoken – including the Prince's true name. She has started to forget. The forgetting is a mercy.
    \item \textbf{The Ninth Coal (The Ninth):} The coal that was never lit; the promise that the Prince will not return. It is warm. It is warm. It is warm.
    \item \textbf{The Silent Smith's Patron (Mykkiel):} The law folded into steel; the covenant that cannot be broken – unless the Prince's shadow crosses it. The shadow is crossing it now.
  \end{itemize}
\end{itemize}

\subsection*{Lore Echoes (Cross‑Expansion Hooks)}
\label{sec:pereshi-lore}

\begin{itemize}
  \item \textbf{Fire Temples and the Covenant (Mykkiel):} Covenant judges recognise the authority of flame‑keepers. A fire‑temple token can be used as a witness in a Covenant court – the flame itself is considered an incorruptible witness. The Prince once stood before a Covenant judge. He lied. The flame flickered, but did not burn. The judge was paid. The judge's ghost now tends a temple in the high passes, still waiting for the truth. The flame flickers when the judge's ghost speaks. It flickers in the same rhythm as the coal.
  \item \textbf{The Ninth Verse and the Ukwe (Sekogo):} The missing verse of the Pereshi epic is said to be carved on the Ukwe's oldest root. The Lethai‑ar know the verse but will not recite it – doing so would call the Hollow. The verse is the Prince's own confession. The Ukwe's root has begun to bleed sap the colour of rust. The sap smells of the Prince's garden.
  \item \textbf{Silent Smith's Swords and the Aeler (Ngomebe):} The silent smiths are said to be descendants of Aeler vent‑priors who stayed in the highlands. Their folding technique is a lost art; the Ngomebe masons would pay a fortune for a working forge. The smiths will only trade for a Prince's relic – and the Prince's shadow is now offering them his own chains. The chains are warm. They pulse.
  \item \textbf{The Buried Archive and the Kalrantha Amulet (Dhahara):} The First Emperor's scrolls contain a description of the amulet's creation. The Veiled Sisters have been searching for the archive for a century. The party may find it first – or find it already looted. The Prince's ghost is also looking. He believes the amulet can restore his name. He is wrong. The amulet can only give him a new one. He does not want a new one.
  \item \textbf{The Prince's Ghost (Red Desert Connection):} Some say the Prince did not die. He walked into the Red Desert and became part of it. His ghost haunts the high passes on stormy nights, asking for a cup of water. Do not give it to him. He will ask for your name next. Lately, he has started asking for a shadow to borrow. He returns them. They are never quite the same.
\end{itemize}

\subsection*{Pereshi Superstitions (burn for +1 die once)}
\begin{itemize}
  \item \textbf{Bread to the Fire:} Crumble a crust into the temple brazier. Ignore the first \emph{Name‑Debt Called} penalty. \hfill `[SUPERSTITION]`
  \item \textbf{Knot the Scroll:} Tie a single ribbon around a book you are returning. Cancel \emph{Forged Scroll} or take –1 die on your next legal roll (your choice). \hfill `[SUPERSTITION]`
  \item \textbf{Name to the Bridge:} Whisper a dead poet's name into the water below. Gain +1 Effect when reciting a verse this scene, but mark Obligation to the Bridge's spirits. \hfill `[SUPERSTITION]`
  \item \textbf{Coal to the Shadow:} Drop a glowing coal into a basin of water while speaking the Prince's forbidden name backwards. The next time the Prince's shadow appears, it will hesitate for one exchange – but it will remember your face. It will remember your face in your dreams. \hfill `[SUPERSTITION]`
\end{itemize}

\subsection*{Thematic SB Spend Table}
\label{sec:pereshi-sb}

\paragraph{Minor Complications (1 SB)}
\begin{itemize}
\item \textbf{Exposure:} Your actions draw unwanted attention from \textbf{fire‑keepers or archivists}.
\item \textbf{Noise:} Sounds of your actions alert nearby \textbf{caravan guards or Tulkani traders}.
\item \textbf{Trace:} Evidence of your passage marks your route for \textbf{Prince's shadow or name‑hunters}.
\item \textbf{Delay:} A brief but meaningful setback costs you \textbf{time or safe passage through a pass}.
\item \textbf{Supply Strain:} Mark +1 segment on a relevant \textbf{resource clock}.
\end{itemize}

\paragraph{Moderate Setbacks (2 SB)}
\begin{itemize}
\item \textbf{Alarm Raised:} \textbf{Local Flame‑Keeper or lord's herald} becomes aware and begins responding.
\item \textbf{Position Lost:} You lose advantageous ground/cover/stealth due to \textbf{ashfall or a bridge recitation failure}.
\item \textbf{Foe Appears:} A \textbf{name‑hungry spirit or a Fhara debt collector} arrives on scene.
\item \textbf{Gear Trouble:} A piece of equipment becomes \textbf{Compromised/Neglected}.
\item \textbf{Lock/Barrier:} A simple obstacle now requires a test to overcome.
\end{itemize}

\paragraph{Serious Trouble (3 SB)}
\begin{itemize}
\item \textbf{Reinforcements:} Additional \textbf{fire‑temple guards, Pereshi riders, or Prince's shadow fragments} arrive.
\item \textbf{Key Gear Breaks:} A crucial tool/weapon becomes temporarily unusable.
\item \textbf{Major Twist:} The situation fundamentally changes – \textbf{name debt called / ninth coal cracks / archive floods}.
\item \textbf{Rail Tick:} Advance a relevant campaign/front clock by 1 segment.
\item \textbf{Condition Applied:} Mark \textbf{Fatigue 1/Harm 1/Condition} appropriate to fiction.
\end{itemize}

\paragraph{Major Turns (4+ SB)}
\begin{itemize}
\item \textbf{Trap Springs:} A prepared danger activates with full effect.
\item \textbf{Authority Arrival:} \textbf{Ninth Flame‑Keeper, Archivist Kaelen, or the Prince's shadow} intervenes.
\item \textbf{Scene Shift:} The environment changes dramatically – \textbf{temple burns / archive collapses / ninth coal shatters}.
\item \textbf{Patron Omen:} Divine/arcane forces take notice – \textbf{omen appears / blessing lost / curse manifests}.
\item \textbf{Narrative Pivot:} The story takes an unexpected turn that reframes objectives.
\end{itemize}

\paragraph{Region-Specific SB Options}
\begin{itemize}
\item \textbf{Pereshi (Fire Law):} Flames change colour, scrolls ignite without heat, a coal pulses out of rhythm.
\item \textbf{Pereshi (Name Law):} A character forgets their own name for a moment, a ghost asks who you are, a written word turns to ash.
\item \textbf{Pereshi (Bridge Law):} A bridge recites a verse from a dead language, a plank snaps at the ninth step, the wind carries a warning.
\end{itemize}

\subsection*{Pursuit Ladder (Near / Mid / Far)}
\begin{itemize}
\item \textbf{Advance/Withdraw (Wits + Survival):} On a hit, shift one band. On a strong hit, also impose \emph{Ash Cloud}.
\item \textbf{Cut the Name (Presence + Sway):} On a hit, freeze range for one exchange; if a rose from the Ruined Garden is in your possession, +1 die.
\item \textbf{Weather the Shadow (Resolve + Spirit):} On a hit in \emph{Prince's Shadow Points}, you pick who is targeted; on a miss, the GM does.
\end{itemize}

\subsection*{Boss Seeds (Ash, Name \& Heart)}
\begin{itemize}
  \item \textbf{The Name‑Eater (Nine Erasures)} – \emph{Phases}: Name Debt Called $\rightarrow$ Orchestrated Forgetting $\rightarrow$ Archive Coup. \emph{Cracks}: return a severed name to its owner, expose a false erasure, survive a night in the Archive of Severed Names.
  \item \textbf{The Ninth Flame‑Keeper (Legend)} – \emph{Phases}: Coal Pulses $\rightarrow$ Shadow Stretches $\rightarrow$ The Heart Cracks. \emph{Cracks}: offer a new name to bind (a willing sacrifice), speak the Prince's true name in the Well of Unspoken Names, or convince the Pereshi that the Prince's debt is finally paid.
  \item \textbf{The Prince's Shadow (Manifestation)} – \emph{Phases}: Whisper $\rightarrow$ Demand $\rightarrow$ Possession. \emph{Cracks}: burn a rose from the Ruined Garden in his presence, speak his true name while holding the ninth coal, or offer him a shadow in exchange for his own – but the shadow must be willing.
\end{itemize}

\begin{tcolorbox}[colback=black!3,colframe=black!40!white,title={Patronage \& Power}]
Among the Pereshi, power flows through fire, ink, and the weight of ancient witness. A lord's authority rests on the age of his library; a flame‑keeper's influence on the purity of her fire. The petty lords are constantly forging genealogies, and the archivists are constantly exposing them. Power is not inherited – it is cited, and it is witnessed by the flame. But now a new power stirs: the Prince's shadow, offering forgotten knowledge, impossible bargains, and the chance to rewrite history. Some Pereshi have begun to listen. The ninth coal is warm. The shadow is patient.

\textbf{For the GM:} Use the \textbf{Prince's Shadow Clock} as the campaign's ticking heart. Every session, describe how the ninth coal feels warmer, how a mirror in a lord's hall has begun to show a second figure, how a scribe's shadow now moves independently. Let the party decide whether to reseal the Prince, bargain with him, or – most dangerously – try to seize his power for themselves. The Pereshi are not a people waiting for a saviour. They are a people waiting for the fire to decide. And the fire is listening. The coal is warm. The shadow is stretching.

Remember the relationships: the Fhara send envoys who cannot meet Pereshi eyes; the Tulkani leave knots at thresholds and wait for the smoke. The lowlands do not see the Pereshi at all. That is how the Pereshi like it. But when the coal cracks, everyone will see – and it will be too late.
\end{tcolorbox}

% Index entries
\index{Pereshi|see{Roof of the Way}}
\index{Theme generator}
\index{Roof of the Way generator}
\index{Places!Pereshi}
\index{People!Pereshi}
\index{Complications!Pereshi}
\index{Rewards!Pereshi}
\index{Flame-Keeper}
\index{Forgotten Prince}
\index{Ninth coal}
\index{Name erasure}
\index{Fhara (fear)}
\index{Tulkani (unease)}
\index{Mount Khvarena}
\index{Archive of Severed Names}
\index{Prince's Shadow clock}
\index{Local Patrons!Pereshi}